% 読み物版の表紙。開いた 1 ページ目として機能する。
%
% **Kindle 端末のライブラリ一覧には出ない。**一覧の表紙は Amazon のサーバーから
% 取りに行く設計で、自分で送った文書には対応レコードがないため、汎用枠に
% タイトル文字だけが出る。ロック画面と Kindle アプリでは出る。
% したがってここは「一覧で目立たせる」ためではなく「開いたときの顔」として作る。
%
% TikZ のノード内で \\ を使うとフォントサイズが 2 行目に引き継がれないので、
% **行ごとに別ノードを置く**。
%
% 引数: {ラベル}{題1行目}{題2行目}{副題}{日付}
% 題が 1 行で足りるときは 3 番目を空にする。3 行要るときはノードを 1 つ足し、
% 以降の y 座標を 16mm ずつ下げる。
% 2026-10-21: 印刷でインクを使いすぎないよう、白地に細い罫線と文字だけの表紙にした。
% 飾りは淡い色の線画（SROC 曲線とフォレストプロット）にとどめる。
\newcommand{\thumbcover}[5]{%
  \thispagestyle{empty}%
  \null
  \begin{tikzpicture}[remember picture, overlay]
    \fill[kkaccent]
      ([yshift=-14mm]current page.north west)
      rectangle ([yshift=-15.2mm]current page.north east);

    \node[anchor=north west, text=kkaccent, inner sep=0pt]
      at ([xshift=12mm, yshift=-28mm]current page.north west)
      {\kkgt{5}\fontsize{11}{11}\selectfont #1};

    \node[anchor=north west, text=kkdark, inner sep=0pt]
      at ([xshift=12mm, yshift=-38mm]current page.north west)
      {\kkgt{7}\fontsize{32}{32}\selectfont #2};
    \node[anchor=north west, text=kkdark, inner sep=0pt]
      at ([xshift=12mm, yshift=-54mm]current page.north west)
      {\kkgt{7}\fontsize{32}{32}\selectfont #3};

    \fill[kkaccent]
      ([xshift=12mm, yshift=-55mm]current page.north west)
      rectangle ([xshift=34mm, yshift=-55.9mm]current page.north west);

    \node[anchor=north west, text=kkdark, inner sep=0pt, align=left,
          font=\kkgt{5}\fontsize{13}{18}\selectfont]
      at ([xshift=12mm, yshift=-62mm]current page.north west)
      {#4};

    % 飾り: 左にフォレストプロット、右に SROC 曲線（淡い線画）
    \begin{scope}[shift={([xshift=14mm, yshift=40mm]current page.south west)}, x=1mm, y=1mm,
                  draw=kksecondary, line width=.6pt]
      \draw[dashed] (28,0) -- (28,62);
      \foreach \y/\c/\w in {58/20/9, 52/24/6, 46/18/12, 40/26/5, 34/22/8, 28/25/7, 22/19/10, 16/23/6} {
        \draw (\c-\w,\y) -- (\c+\w,\y);
        \fill[kksecondary] (\c-1.2,\y-1.2) rectangle (\c+1.2,\y+1.2);
      }
      \fill[kkaccent!60] (16,6) -- (22,8) -- (28,6) -- (22,4) -- cycle;
    \end{scope}
    \begin{scope}[shift={([xshift=-74mm, yshift=40mm]current page.south east)}, x=1mm, y=1mm,
                  line width=.6pt]
      \draw[kksecondary] (0,0) rectangle (60,60);
      \draw[kksecondary, dotted] (0,0) -- (60,60);
      \draw[kkaccent!70, line width=1pt] plot[domain=0:1, samples=60, smooth] ({60*\x}, {60*(\x)^(0.28)});
      \foreach \px/\py in {14/46, 18/52, 9/40, 22/50, 12/44, 26/55, 16/38, 30/54} {
        \draw[kksecondary, fill=white] (\px,\py) circle (1.5);
      }
      \draw[kkaccent, fill=kkwarm!70] (16,47) -- (17.8,48.8) -- (16,50.6) -- (14.2,48.8) -- cycle;
    \end{scope}

    \fill[kkaccent]
      ([yshift=22mm]current page.south west)
      rectangle ([yshift=22.6mm]current page.south east);
    \node[anchor=south west, text=kkdark, inner sep=0pt]
      at ([xshift=12mm, yshift=11mm]current page.south west)
      {\kkgt{3}\fontsize{9}{9}\selectfont #5};
  \end{tikzpicture}%
  \clearpage}
